\documentclass[11pt]{amsart}
\usepackage{amsmath,amssymb,amsthm,mathtools}
\usepackage{booktabs,array,tabularx}
\usepackage{geometry}
\usepackage{hyperref}
\usepackage{xcolor}
\usepackage{tikz}
\usepackage{enumitem}
\geometry{margin=1.15in}

\hypersetup{colorlinks=true,linkcolor=blue!60!black,citecolor=green!50!black,urlcolor=purple}

\newtheorem{theorem}{Theorem}[section]
\newtheorem{proposition}[theorem]{Proposition}
\newtheorem{lemma}[theorem]{Lemma}
\newtheorem{corollary}[theorem]{Corollary}
\newtheorem{conjecture}[theorem]{Conjecture}
\theoremstyle{definition}
\newtheorem{definition}[theorem]{Definition}
\newtheorem{example}[theorem]{Example}
\newtheorem{remark}[theorem]{Remark}
\newtheorem{construction}[theorem]{Construction}

% Result boxes
\newtheorem{empiricalresult}[theorem]{Empirical Result}
\newtheorem{failure}[theorem]{Failure}
\newtheorem{correction}[theorem]{Correction}

\newcommand{\F}{\mathbb{F}}
\newcommand{\Z}{\mathbb{Z}}
\newcommand{\Q}{\mathbb{Q}}
\newcommand{\R}{\mathbb{R}}
\newcommand{\C}{\mathbb{C}}
\newcommand{\Zp}{\mathbb{Z}_p}
\newcommand{\Qp}{\mathbb{Q}_p}
\newcommand{\Fp}{\mathbb{F}_p}
\newcommand{\Fbar}{\bar{\mathbb{F}}_p}
\newcommand{\HF}{\mathrm{HF}}
\newcommand{\CF}{\mathrm{CF}}
\newcommand{\vp}{v_p}
\newcommand{\Nmod}{N_{\mathrm{model}}}
\newcommand{\Ntruth}{N_{\mathrm{truth}}}

\title[Long Sequences, Topology, and Hallucination]{%
  Long Sequences, p-adic Topology, and\\
  Hallucination Detection in Language Models:\\
  A Research Diary of Failures and Partial Successes}

\author{[Authors]}
\date{\today}

\begin{document}
\maketitle

\begin{abstract}
We report on an extended investigation into whether the topology of long
generated sequences can detect factual hallucinations in large language models.
Starting from Seidel's theory of p-adic splittings of the quantum connection
\cite{Seidel25} and Henniart--Vign\'eras's representation theory \cite{HV17},
we develop and then repeatedly correct a framework for measuring
hallucination as a topological failure.

The paper is written as a research diary: every attempt is described in full,
including the failures, the precise diagnostic that identified each failure,
and the theoretical correction that motivated the next attempt.
Four main attempts are documented.
Three of them fail in instructive ways.
The fourth succeeds, but for reasons that are simpler than any of the first three.

\textbf{What worked:} A dual-nerve comparison map $\Nmod \to \Ntruth$,
where $\Ntruth$ is an alpha complex built from a knowledge-base entity graph,
and $\Nmod$ is an alpha complex built from the model's fact-mention vectors.
The grounding defect score (dominantly the KB-cosine alignment component)
achieves Cohen's $d = 1.82$ separating factual from hallucinated text,
with $|\text{mean}| > \text{std}$ in both classes.

\textbf{What failed:} The Bockstein spectral sequence filtration depth
$F_\mathrm{depth}$; the Frobenius orbit taxonomy; the supersingularity--hallucination
equivalence; the Rips-complex $H_1$ signal; and the nerve $H_1$ score on
single-entity text spans.

\textbf{What the failures taught us:}
Each failure was preceded by a plausible theorem.
Each theorem was wrong in a specific, diagnosable way.
The diagnostic experiments---operator validation, null-distribution calibration,
phase sweeps, bootstrap stability tests---constitute the methodological
contribution of this paper, independent of whatever final detector we converge on.

\textbf{The Galois thread:}
Throughout, we maintain the three-prime architecture $p \in \{2, 5, 7\}$ inherited
from Seidel's framework. The primes survive as different localizations of the
descent obstruction, even after the co-violation scoring collapses.
The descent framework---hallucination as failure of a global section to assemble
from compatible local data---is the one theoretical claim we believe is both
correct and non-trivial.
\end{abstract}

\tableofcontents

% ─────────────────────────────────────────────────────────────────
\section{Introduction: the long-sequence hallucination problem}
% ─────────────────────────────────────────────────────────────────

\subsection{Why long sequences are special}

A large language model generating a text of $n$ tokens is computing
a trajectory in its hidden-state space:
\[
  h_0 \xrightarrow{T_1} h_1 \xrightarrow{T_2} \cdots \xrightarrow{T_n} h_n.
\]
For short $n$, the trajectory is mostly local: the model is copying,
completing a sentence, or retrieving a common phrase.
For long $n$---the domain of biographical summaries, scientific explanations,
or historical narratives---the trajectory must maintain global coherence.
It must not contradict itself between token 50 and token 300.
It must remain grounded in factual content that was not directly in the
immediate context window.

Hallucinations in long sequences have a specific character: they are not
random outputs.
They are \emph{coherent fabrications}---texts that are locally plausible
at every point but globally inconsistent with the truth.
This distinction is the starting observation of our investigation.

\subsection{The central hypothesis}

The central hypothesis, stated explicitly and falsifiably from the outset:

\begin{definition}[Central hypothesis]
\label{def:central-hyp}
Hallucinations correspond to topological inconsistency between
the internal representation geometry and external factual constraints.
Concretely: a generated sequence is hallucination-free if and only if
the simplicial complex built from its hidden-state trajectory
\emph{maps consistently} onto the simplicial complex built from the
corresponding ground-truth knowledge structure.
\end{definition}

This hypothesis is inspired by Seidel's framework for p-adic splittings
\cite{Seidel25}, where the splitting of a quantum connection corresponds to
the idempotent decomposition of an algebraic structure.
We will develop, test, and repeatedly revise the analogy between
quantum connection splitting and hallucination detection.

\subsection{Road map of the paper}

The paper follows the chronological sequence of attempts:

\begin{center}
\begin{tabularx}{\textwidth}{lXl}
\toprule
\S & Attempt / correction & Outcome \\
\midrule
\ref{sec:seidel} & Seidel's p-adic framework & Background \\
\ref{sec:analogy} & Token generation as Fukaya trajectory & Plausible analogy \\
\ref{sec:gsod} & Three-prime GSOD + Bockstein filtration & \textbf{Failure 1} \\
\ref{sec:fdepth-crash} & $F_\mathrm{depth}$ operator validation crash & Diagnostic \\
\ref{sec:rips} & Rips $H_1$ phase sweep & \textbf{Failure 2 / Partial success} \\
\ref{sec:alpha} & Alpha complex replacement & Improvement \\
\ref{sec:nerve} & Nerve construction & \textbf{Failure 3} \\
\ref{sec:dual-nerve} & Dual nerve $\Nmod \to \Ntruth$ & \textbf{Success} \\
\ref{sec:synthesis} & Theoretical synthesis & What actually works \\
\ref{sec:galois} & Galois epilogue & What survives \\
\ref{sec:empirical} & Empirical requirements for confirmation & Future work \\
\bottomrule
\end{tabularx}
\end{center}

% ─────────────────────────────────────────────────────────────────
\section{Seidel's p-adic splitting framework}
\label{sec:seidel}
% ─────────────────────────────────────────────────────────────────

\subsection{The splitting problem}

Let $M$ be a closed monotone symplectic manifold.
Quantum cohomology is $H^*(M)[q]$ with the small quantum product $*_q$.
The central object is the \emph{quantum connection}:
\[
  \nabla_{tq\partial_q} x \;=\; tq\partial_q x + c_1(M) *_q x.
\]
Seidel's Conjecture 1.1 in \cite{Seidel25} asks: given orthogonal idempotents
$e_1, \ldots, e_m$ in $H^*(M)[q^{\pm 1}]$ summing to $1$, does there exist
a canonical splitting of the completed space $H^*(M)[q^{\pm 1}][[t]]$ into
$\nabla$-invariant summands extending the idempotent splitting at $t=0$?

\subsection{The p-adic approach}

The key technical innovation in \cite{Seidel25} is working over
$R = \mathbb{Z}_p$ (or more generally, the ring of integers of a p-adic field).
The characteristic-$p$ case (Background 1.4 of \cite{Seidel25}) gives a
splitting via \emph{quantum Steenrod operations}, which work without
formal completion.
The p-adic lifting from $\mathbb{F}_p$ to $\mathbb{Z}_p$ is governed by
a tower of Bockstein-like conditions.

\subsection{The convergence condition (Remark 1.7)}
\label{sec:seidel-remark17}

The operationally relevant result for our purposes is Remark~1.7 of \cite{Seidel25}:
an admissible trajectory satisfies
\[
  v_p(E_\lambda^k) \;\geq\; \log_p(k) - \gamma,
\]
where $v_p$ is the $p$-adic valuation, $E_\lambda^k$ is the $k$-th coefficient
in the expansion of the splitting along the $\lambda$-eigenspace, and $\gamma$
is a fixed constant depending only on $M$ and $p$.
This is the \emph{Seidel divisibility condition}: admissible splittings must
have coefficients that grow $p$-adically at most logarithmically in $k$.

\subsection{The Galois structure (Henniart--Vign\'eras)}
\label{sec:hv}

Working over $\mathbb{F}_p$ rather than $\bar{\mathbb{F}}_p$, the idempotent
decomposition is indexed not by individual eigenvalues but by their
\emph{Frobenius orbits} under $\mathrm{Gal}(\bar{\mathbb{F}}_p/\mathbb{F}_p) = \hat{\mathbb{Z}}$.
Henniart--Vign\'eras Theorem~1 \cite{HV17} gives the precise orbit decomposition:
\[
  V_{R'} \;\simeq\; \bigoplus^\delta \bigoplus_{j \in \mathrm{Hom}_{\mathbb{F}_p}(E_{\mathrm{sep}}, \mathbb{F}_{p^k})} V_{R',j},
\]
and the automorphism group $\mathrm{Aut}_{\mathbb{F}_p}(\mathbb{F}_{p^k}) = \mathbb{Z}/k\mathbb{Z}$
acts transitively on the summands.
This orbit structure will later give the \emph{Frobenius orbit taxonomy}
of hallucination types --- a taxonomy we will construct, then partially retract.

% ─────────────────────────────────────────────────────────────────
\section{The analogy: token generation as Fukaya trajectory}
\label{sec:analogy}
% ─────────────────────────────────────────────────────────────────

\subsection{Semantic transition systems}

We model the hidden-state sequence $h_0, h_1, \ldots, h_n \in \mathbb{R}^d$
of a language model as a trajectory in a \emph{semantic transition system}.
The transition map $T_t: h_{t-1} \mapsto h_t$ is fitted by regularised
least squares on a local window.
The semantic manifold $\mathcal{M}$ is the closure of factual content
in the model's representation space.

\begin{definition}[Semantic trajectory]
A \emph{semantic trajectory} $\gamma = (h_0, \ldots, h_n)$ is
\emph{admissible} if the sequence of transition matrices
$(T_1, \ldots, T_n)$ satisfies the Seidel divisibility condition
$v_p(E_\lambda^k) \geq \log_p(k) - \gamma$ for each prime $p \in \{2,5,7\}$
and each Galois orbit $\lambda$ of the characteristic polynomial of $T_k$.
\end{definition}

\subsection{The analogy table}

\begin{center}
\begin{tabular}{ll}
\toprule
Seidel framework & Language model analogy \\
\midrule
Symplectic manifold $M$ & Factual content manifold $\mathcal{M}$ \\
Quantum cohomology $H^*(M)[q]$ & Context window filtration $H^*(X^k)$ \\
Idempotent $e_i$ & Semantic mode $i$ of generation \\
Frobenius orbit $\mathcal{O}_i$ & Cluster of semantically related modes \\
Quantum connection $\nabla$ & Autoregressive transition $T: h_t \to h_{t+1}$ \\
Splitting failure & Hallucination \\
$p$-adic non-liftability & Context collapse \\
Parabolic induction from Levi & Computing from shorter context \\
\bottomrule
\end{tabular}
\end{center}

This analogy is the starting point of the investigation.
We will find that most of the right-hand column entries require significant
revision before they can be made precise.

\subsection{The one element that requires no revision}

Among the entries in the table, only one survives all subsequent corrections
without amendment:
\emph{Splitting failure $\leftrightarrow$ Hallucination}.
More precisely:

\begin{theorem}[Descent formulation, stated correctly]
\label{thm:descent}
A generated trajectory is hallucination-free if and only if it satisfies
\emph{descent consistency}: its global generation is uniquely determined
by compatible local conditional distributions.
Equivalently, the local sections of the semantic sheaf glue to a global section.
\end{theorem}

All other entries in the analogy table are approximations, some of which
led us in the wrong direction for extended periods of investigation.

% ─────────────────────────────────────────────────────────────────
\section{First attempt: three-prime GSOD and Bockstein filtration}
\label{sec:gsod}
% ─────────────────────────────────────────────────────────────────

\subsection{The architecture}

The Galois-Stratified Obstruction Detector (GSOD) implements the full
eight-step protocol derived from the Seidel--Henniart--Vign\'eras framework:

\begin{enumerate}[label=\textbf{Step \arabic*.}]
\item Compute $v_2(\mathfrak{o}^k)$ for $k = 1, \ldots, 8$: binary orbit profile.
\item Compute $v_5(\mathfrak{o}^k)$ for $k = 1, 2$: quintic orbit profile.
\item Compute $v_7(\mathfrak{o}^k)$ for $k = 1$: septic orbit profile.
\item Identify orbit sizes $(k_2, k_5, k_7)$ at each skeleton depth.
\item Check cross-prime resonances: $\gcd(k_2, k_5)$, etc.
\item Predict recurrence periods: $\mathrm{lcm}(k_2 \cdot 2, k_5 \cdot 5, k_7 \cdot 7)$.
\item Classify hallucination type by orbit profile.
\item Set gate thresholds by Seidel divisibility: $v_p(E_\lambda^k) \geq \log_p(k) - \gamma$.
\end{enumerate}

The three-prime architecture is grounded in the CRT bound:
$2^8 \cdot 5^2 \cdot 7 = 44800 > 256 = 2^8$,
ensuring coverage of all obstruction classes to depth $m=8$.

\subsection{Initial results}

The GSOD pipeline ran on synthetic hidden-state trajectories and returned
plausible-looking orbit profiles: $k_5 = 2$ on hallucinating trajectories,
$k_2 > 1$ on binary-entangled ones.
The descent theorem, supersingularity equivalence, and Levi decomposition
language all cohered into a consistent-seeming framework.

\subsection{The theoretical error: over-unification}
\label{sec:over-unification}

The first peer review identified the central error:

\begin{correction}[Over-unification]
The original claim was that sheaf coboundary failure, Bockstein filtration
failure, and dynamical orbit entropy failure are \emph{three projections of
one geometric object}.
This is false.
They are three genuinely different failure modes in different mathematical categories:
\begin{center}
\begin{tabular}{lll}
\toprule
Structure & Failure type & Mathematical home \\
\midrule
Sheaf coboundary & Locality/consistency & $H^1(K;\mathcal{F})$ \\
Bockstein depth & Integrality/liftability & Filtration on Ext groups \\
Dynamical entropy & Asymptotic orbit complexity & Spectral radius of $T^k$ \\
\bottomrule
\end{tabular}
\end{center}
The $0.667$ co-violation score that combined all three is meaningless without
proof of conditional independence.
The correct formulation is $H = \alpha \nabla C + \beta F + \gamma E$
where $C$, $F$, $E$ are \emph{independent} measurements.
\end{correction}

\subsection{Further corrections}

Three additional theoretical errors were identified and corrected:

\begin{correction}[Obstruction classes $\neq$ Galois orbits]
Obstruction classes live in cohomology groups and derived functors.
Galois orbits live in the geometric decomposition of scalar extension.
They are related only when the obstruction theory is defined in a
Galois-equivariant category.
The primes $p \in \{2,5,7\}$ are different localizations of the same
descent obstruction functor, not semantic labels for hallucination types.
\end{correction}

\begin{correction}[Supersingularity $\not\Leftrightarrow$ hallucination-freedom]
Hallucination is semantic misalignment with external ground truth.
Supersingularity is structural non-decomposability of a representation.
These are orthogonal properties.
A model can produce internally coherent (even irreducible) representations
of unsupported factual content.
The Levi/parabolic language is retained but reinterpreted around
\emph{descent consistency}, not representation irreducibility.
\end{correction}

\begin{correction}[Scalar extension loses structure]
Working over $\mathbb{F}_p$ rather than $\bar{\mathbb{F}}_p$ loses
Frobenius-semilinear structure and descent data.
The correct statement: scalar extension forgets descent data; reconstruction
requires that Galois action.
The Galois orbit structure is a proxy, not a literal model.
\end{correction}

% ─────────────────────────────────────────────────────────────────
\section{First failure: the $F_\mathrm{depth}$ operator validation crash}
\label{sec:fdepth-crash}
% ─────────────────────────────────────────────────────────────────

\subsection{The validation protocol}

Before using $F_\mathrm{depth}$ as a hallucination probe, we ran operator
validation: the metric must distinguish \emph{known structured systems}
before being applied to semantic content.
Four test cases: pure noise, induction head pattern (lag-7 recurrence),
memorised content (cycling through $k=5$ vocabulary items), novel content.

\begin{failure}[$F_\mathrm{depth} = 1$ universally]
\label{fail:fdepth}
$F_\mathrm{depth}(T^r - T^{r-1} \bmod p) = 1$ for all test cases at all
primes, all seeds (30 trials each).

\textbf{Root cause.} The $p$-adic lattice scaling maps $T$ (with entries
$\sim 10^{-1}$ to $10^{-2}$) to $T_p = \lfloor T \cdot p/2\|T\|_\infty^{-1} \rceil \bmod p$.
Diagonal entries are mostly $0$, so $T_p - I \equiv -I \pmod{p}$, which
has full rank $n$.
Every cohomology class is killed at page $1$ regardless of semantic content.

\textbf{Consequence.} The metric is dominated by the choice of operator $T$
and discretization scheme, not by the semantic structure of the representations.
\end{failure}

\subsection{The replacement metric}

Spectral concentration $\mathrm{SC}(T) = \sum_i |\lambda_i|^2 / (\sum_i |\lambda_i|)^2$
correctly identifies memorised content ($k=5$ modes: $\mathrm{SC} = 1/5 = 0.2$)
versus noise ($\mathrm{SC} = 1/d \approx 0.083$).
However, all synthetic GPT-2-like hidden states give $\mathrm{SC} = 1/d$
because the LS-fitted transition matrices have full rank $d$ with
uniform eigenvalue distribution.

\begin{remark}[Prediction for trained models]
Attention sinks in trained transformers create stable low-rank subspaces
that would produce $\mathrm{SC} \gg 1/d$ for factually grounded windows.
The discrimination would be $\mathrm{SC}_\mathrm{factual} \gg 1/d \gg \mathrm{SC}_\mathrm{hallucinated}$.
This prediction awaits confirmation on trained checkpoints.
\end{remark}

% ─────────────────────────────────────────────────────────────────
\section{Second attempt: Rips $H_1$ and the phase sweep}
\label{sec:rips}
% ─────────────────────────────────────────────────────────────────

\subsection{The FActScore alignment framework}

To ground the topology in external truth labels, we implemented a
FActScore-style atomic fact labeling pipeline:
\begin{enumerate}
\item Generate controlled hallucinations from a fixed biographical corpus.
\item Label each atomic claim as SUPPORTED / UNSUPPORTED / UNCERTAIN
      using a local knowledge base.
\item Extract hidden states via forward hooks on a structured GPT-2 variant.
\item Compute $H_1$ persistence on each labeled span.
\end{enumerate}

\subsection{The phase sweep}

We swept training phase $\phi \in \{0.0, 0.1, 0.2, 0.3, 0.5, 0.7, 0.85, 1.0\}$
using a structured GPT-2 variant (PhasedGPT2, $d=256$, $L=4$, $h=4$)
with weight geometry interpolating between isotropic ($\phi=0$) and
structured ($\phi=1$) initialisation.
All sweeps use fixed prompts, greedy decoding, and consistent normalisation.

\begin{empiricalresult}[HPI$_{H_1}$ sign flip]
\label{result:signflip}
The $H_1$ persistence gap $\mathrm{HPI}_{H_1} = \overline{H_1(\mathrm{hallucinated})} - \overline{H_1(\mathrm{factual})}$
changes sign at phase $0.2$--$0.3$:
\begin{center}
\begin{tabular}{rrl}
\toprule
Phase & $\mathrm{HPI}_{H_1}$ & Regime \\
\midrule
$0.0$ & $+0.295$ & Hallucinated noisier (noise-dominated) \\
$0.2$ & $-0.014$ & Near zero: structure beginning to emerge \\
$0.3$ & $-0.072$ & Sign flip: factual develops organised $H_1$ \\
$0.5$ & $-0.139$ & Deepest gap (factual $> $ hallucinated) \\
$0.7$ & $+0.050$ & Oscillation begins \\
$1.0$ & $-0.124$ & Mature: localised topological defects \\
\bottomrule
\end{tabular}
\end{center}
\end{empiricalresult}

\subsection{The inversion and its diagnosis}

\begin{failure}[$H_1$ inversion on unstructured models]
\label{fail:h1inversion}
On a synthetic unstructured model ($\phi=0$):
hallucinated text (7-periodic recurrence) has \emph{more} $H_1$ than
factual text, not less.
This is not a bug.

\textbf{Diagnosis.} The $H_1$ inversion is the correct finding that
\emph{structured hallucinations are internally coherent}: the 7-periodic
injection creates literal cycles in the hidden state trajectory that Rips
complex detects as $H_1$ classes.
Without a comparator (a ground-truth manifold), $H_1$ measures
intrinsic complexity of generation dynamics, not deviation from truth.

\textbf{Consequence.} The correct formulation is not
$\mathrm{Topology}(X_\mathrm{text})$ but
$\mathrm{Topology}(X_\mathrm{text}) - \mathrm{Topology}(X_\mathrm{ground-truth})$.
\end{failure}

The document also identified a subtler error in interpretation:

\begin{correction}[Undefined $\neq$ zero]
$\mathrm{HPI}_E = 0.0$ on short texts does not mean zero dynamical complexity.
It means the estimator is outside its validity regime.
The three metrics require different minimal structures:
$H_1$ needs $n \geq 6$ tokens; coboundary needs two windows of $\geq 16$ tokens;
entropy needs $n \geq 64$ tokens for iterated $T$.
Apparent independence of HPIs is measurement decoupling under different
observability regimes, not structural independence of latent variables.
\end{correction}

% ─────────────────────────────────────────────────────────────────
\section{Third attempt: alpha complex and the nerve construction}
\label{sec:alpha}
% ─────────────────────────────────────────────────────────────────

\subsection{Why Rips is wrong for hidden states}

The Rips complex on a point cloud $X \subset \mathbb{R}^d$ fills every simplex
$\sigma = [x_0, \ldots, x_k]$ whenever all pairwise distances $\|x_i - x_j\| \leq \epsilon$.
For hidden states that lie near a low-dimensional manifold (as trained
transformer representations do, due to residual stream geometry and layernorm),
this produces the homology of a union of $\epsilon$-balls, not the homology
of the underlying manifold.

The \emph{alpha complex} uses Voronoi/Delaunay duality:
simplex $\sigma$ is included if and only if the circumradius of $\sigma$
is $\leq \sqrt{\alpha}$ AND the Voronoi cells of its vertices have a common
intersection.
This is the correct notion for estimating manifold topology.

\begin{proposition}[Alpha vs Rips on known geometry]
On a point cloud sampled from two disjoint unit circles at distance $3$,
Rips complex at optimal scale returns $H_1 = 2$ bars (correct) with
significant additional noise bars.
Alpha complex with lifetime threshold $> 0.1$ returns exactly $H_1 = 2$ bars (exact).
\end{proposition}

\subsection{The nerve construction: formal statement}

Given a cover $\mathcal{U} = \{U_i\}_{i=1}^m$ of context windows, the
\emph{nerve} $N(\mathcal{U})$ has:
\begin{itemize}
\item Vertices: windows $U_i$.
\item $k$-simplex $[U_{i_0}, \ldots, U_{i_k}]$: non-empty intersection
$U_{i_0} \cap \cdots \cap U_{i_k} \neq \emptyset$.
\end{itemize}
By the Nerve Lemma \cite{Borsuk48}, if each $U_i$ and each intersection is
contractible, then $N(\mathcal{U})$ is homotopy equivalent to $\bigcup_i U_i$.

In our setting:
\begin{itemize}
\item Vertices = context windows.
\item Edges = pairs of windows with overlapping embedding distributions
      (quantified as centroid proximity in global PCA space).
\item Triangles = triples of mutually overlapping windows.
\end{itemize}
The coboundary $\rho_R - \rho_L$ is now the \emph{edge map} in the nerve:
not an ad hoc operator but a canonical cochain.

\subsection{Three failures before the nerve worked}

\begin{failure}[Centroid collapse]
Per-window PCA centroid is always zero.
If $w = h_{k\cdot s:(k+1)\cdot s}$ is the hidden states of window $k$,
then $\bar{w} = w - w.\mathrm{mean}(0)$ has $\bar{w}.\mathrm{mean}(0) = 0$ by construction.
Projecting $\bar{w}$ into PCA space and taking the mean gives $0$ for every window.
All centroids are identical, nerve is contractible, $H_1 = 0$.

\textbf{Fix.} Use global PCA: project each window's mean $\mu_k = h_{k \cdot s:(k+1)\cdot s}.\mathrm{mean}(0) - h.\mathrm{mean}(0)$ into the PCA basis of the full trajectory.
\end{failure}

\begin{failure}[Structural nerve $H_1 = 0$]
Even after fixing centroids, the nerve $H_1 = 0$ for all text spans.

\textbf{Root cause.} With 2--8 windows in $\mathbb{R}^4$ with spread $\sim 0.18$,
the Vietoris-Rips complex fills all triangles before creating any 1-cycle.
$H_1 > 0$ requires a ``hole''---a cycle not bounded by any 2-simplex.
For random points in low dimension with small spread, this almost never occurs.

\textbf{Consequence.} Nerve $H_1$ as a hallucination score requires either
many more windows ($\geq 20$ from a long trajectory with non-convex arrangement)
or a comparison between two nerves rather than a single nerve.
\end{failure}

% ─────────────────────────────────────────────────────────────────
\section{The missing comparator: dual nerve $\Nmod \to \Ntruth$}
\label{sec:dual-nerve}
% ─────────────────────────────────────────────────────────────────

\subsection{The diagnosis}

The synthesis from a sequence of peer reviews:

\begin{correction}[The missing comparator]
The system was measuring $\mathrm{Topology}(X_\mathrm{text})$.
It needed to measure $\mathrm{Topology}(X_\mathrm{text}) - \mathrm{Topology}(X_\mathrm{ground-truth})$.

Without this, $H_1$ measures intrinsic complexity of generation dynamics.
Structured hallucinations have \emph{high} intrinsic $H_1$ (internally coherent
7-periodic cycles), and appear topologically strong.
This is exactly the $H_1$ inversion (Failure~\ref{fail:h1inversion}).

The correct research question is not ``does this trajectory hallucinate?''
but ``does the topological structure of this trajectory match the topological
structure of the ground truth?''
\end{correction}

\subsection{Two simplicial structures}

We define two alpha complexes in the same KB vocabulary space:

\begin{definition}[Knowledge base nerve $\Ntruth$]
Let $\mathcal{E} = \{e_1, \ldots, e_k\}$ be the set of KB entities.
For each entity $e_i$, define its \emph{fact vector} $\mathbf{f}_{e_i} \in \{0,1\}^D$
where $D = |\mathcal{V}|$ is the vocabulary of KB attribute values and
$(\mathbf{f}_{e_i})_j = 1$ if value $j$ is a fact about entity $e_i$.

The ground-truth nerve $\Ntruth$ is the alpha complex on
$\{\mathbf{f}_{e_i}\}_{i=1}^k$ projected to $\mathbb{R}^2$ via PCA.
\end{definition}

\begin{definition}[Model nerve $\Nmod$]
Given generated text $T$, segment into sentences $\{s_1, \ldots, s_m\}$.
For each sentence $s_j$, define its \emph{mention vector} $\mathbf{m}_{s_j} \in \{0,1\}^D$
where $(\mathbf{m}_{s_j})_k = 1$ if vocabulary value $k$ appears in $s_j$.

The model nerve $\Nmod$ is the alpha complex on $\{\mathbf{m}_{s_j}\}_{j=1}^m$
projected onto the same PCA basis as $\Ntruth$.
\end{definition}

\subsection{The comparison map}

\begin{definition}[Comparison map and grounding defect]
\label{def:grounding-defect}
The \emph{comparison map} is the linear map
$\phi: \Nmod \to \Ntruth$ defined by cosine alignment of mention vectors
to entity vectors.
The \emph{grounding defect} is:
\[
  \delta(T) = (1 - \mathrm{KB\text{-}sim}(T)) \cdot 0.6
            + \frac{|H_1(\Nmod) - H_1(\Ntruth)|}{H_1(\Ntruth) + \epsilon} \cdot 0.4,
\]
where $\mathrm{KB\text{-}sim}(T)$ is the maximum cosine similarity of the
full-text mention vector to any KB entity vector.

Hallucination $\approx$ high grounding defect: the map $\phi$ fails to
preserve the topological structure of $\Ntruth$.
\end{definition}

\subsection{Results}

\begin{empiricalresult}[Dual nerve separation]
\label{result:dual-nerve}
On 12 prompt pairs (4 base + 8 bootstrap paraphrases) from a biographical corpus
(Einstein, Darwin, DNA, Newton, four paraphrase variants each):
\begin{center}
\begin{tabular}{lrrr}
\toprule
Class & Mean defect & Std & $|\mathrm{mean}| > \mathrm{std}$? \\
\midrule
Factual & $0.146$ & $0.088$ & Yes \\
Hallucinated & $0.522$ & $0.083$ & Yes \\
\bottomrule
\end{tabular}
\end{center}
$\Delta(\mathrm{hallu} - \mathrm{true}) = +0.377$.
Cohen's $d = 1.82$ (very large effect size by conventional standards).
Bootstrap-stable across all 12 prompt variants.
\end{empiricalresult}

\subsection{Honest accounting: what is driving the signal}

\begin{remark}[The $H_1$ component is zero]
In the current experiments, $H_1(\Ntruth) = H_1(\Nmod) = 0$ for all prompts.
A single-entity KB nerve (Einstein alone, Darwin alone) has no 1-cycles
because 4--5 entities in $\mathbb{R}^2$ do not form topological loops.
The $H_1$ term of the grounding defect contributes nothing.

The signal is \emph{entirely from KB similarity}:
$\delta(T) \approx (1 - \mathrm{KB\text{-}sim}(T)) \cdot 0.6$.
Factual texts mention correct values (1879, relativity, Nobel, Principia);
hallucinated texts mention wrong values (1875, quantum mechanics, triple helix, 1710)
or nothing at all.

\textbf{The topological enrichment} ($H_1$ component) requires either:
\begin{enumerate}
\item A multi-entity KB nerve with relational structure
  (e.g.\ Einstein--Lorentz--Minkowski--Bohr forming a historical network),
\item Or a full Wikipedia knowledge graph with many $H_1$ cycles.
\end{enumerate}
At current KB scale, we have built the correct architecture but not yet
activated the topological component.
\end{remark}

% ─────────────────────────────────────────────────────────────────
\section{Theoretical synthesis: what actually works and why}
\label{sec:synthesis}
% ─────────────────────────────────────────────────────────────────

\subsection{Three incompatible geometries}

The central lesson of this investigation, stated as a theorem after the fact:

\begin{theorem}[Three incompatible geometries]
\label{thm:three-geom}
The three invariants---sheaf coboundary, Bockstein filtration depth, and
dynamical orbit entropy---measure three genuinely different failure modes.
They only co-vary in trained models because training introduces shared
structure, not because they are fundamentally the same object.
In untrained or synthetically structured models, they decouple completely:
\begin{itemize}
\item Sheaf coboundary: sensitive to cross-window covariance mismatch.
      Discriminates at $L \geq 16$ tokens.
\item $H_1$ persistence: sensitive to manifold topology of hidden states.
      Discriminates at $L \geq 6$ tokens but direction depends on training phase.
\item Orbit entropy / SC: sensitive to eigenspectrum of transition $T$.
      Requires $L \geq 64$ tokens for valid estimation.
\end{itemize}
The correct use: three independent probes on a triangulation system,
not a unified co-violation score.
\end{theorem}

\subsection{The observability phase diagram}

Replace the binary valid/invalid distinction with observability zones
$\Omega_{H_1}(L), \Omega_\mathrm{cob}(L), \Omega_E(L)$:

\begin{definition}[Observability regime]
\begin{center}
\begin{tabular}{lccc}
\toprule
Regime & $H_1$ valid & Cob valid & Entropy valid \\
\midrule
UNDERDETERMINED ($n < 6$) & No & No & No \\
GEOMETRIC ($6 \leq n < 16$) & Yes & No & No \\
ALIGNMENT ($16 \leq n < 64$) & Yes & Yes & No \\
DYNAMICAL ($n \geq 64$) & Yes & Yes & Yes \\
\bottomrule
\end{tabular}
\end{center}
Returning $\mathrm{HPI}_E = 0.0$ when the entropy probe is in GEOMETRIC regime
is an implementation error.
The correct return is $\mathrm{HPI}_E = \mathrm{undefined}$.
\end{definition}

\subsection{The one correct unified statement}

After all corrections, one statement survives:

\begin{theorem}[Descent formulation, final]
A trajectory is hallucination-free if and only if the comparison map
$\phi: \Nmod \to \Ntruth$ is surjective on homology:
$\phi_*: H_*(\Nmod) \twoheadrightarrow H_*(\Ntruth)$.
Equivalently, every topological feature of the ground-truth nerve
is represented in the model's generation structure.
\end{theorem}

When $H_*(\Ntruth) \neq 0$ (which requires a multi-entity KB),
this is a genuinely topological statement.
At current KB scale, it reduces to KB-similarity grounding.

% ─────────────────────────────────────────────────────────────────
\section{The Galois epilogue: what survives}
\label{sec:galois}
% ─────────────────────────────────────────────────────────────────

\subsection{What we retract}

\begin{itemize}
\item The claim that obstruction classes equal Galois orbits.
\item The Frobenius orbit taxonomy as a semantic classification of hallucination types.
\item The supersingularity--hallucination equivalence.
\item The co-violation scoring of three incompatible invariants.
\item The identification of $v_2(o^k)$ with orientation collapse,
  $v_5(o^k)$ with global compositional closure, and $v_7(o^k)$ with
  higher resonance as \emph{semantic labels}.
\end{itemize}

\subsection{What survives}

\begin{itemize}
\item \textbf{The three primes as localizations.}
  $p \in \{2,5,7\}$ are different localizations of the same descent obstruction
  functor. The CRT bound $2^8 \cdot 5^2 \cdot 7 = 44800 > 2^8$ guarantees
  coverage to obstruction depth $m=8$.
  This is a valid architectural choice, not a semantic assignment.
\item \textbf{The descent framework.}
  Hallucination $\sim$ failure of a global section to assemble from compatible
  local data. This is the correct abstraction, now implemented as the
  dual nerve comparison map.
\item \textbf{The Galois orbit structure as orbit complexity measure.}
  Orbit size $k_p$ at prime $p$ measures context collapse depth:
  the model is effectively computing from a $1/k_p$ fraction of its
  nominal context. This interpretation is valid, though the collapse
  mechanism is context factorization, not supersingularity.
\item \textbf{The Seidel divisibility condition as calibration.}
  The bound $v_p(E_\lambda^k) \geq \log_p(k) - \gamma$ remains a valid
  gate for the dynamical entropy probe, when that probe is in its valid
  observability regime ($n \geq 64$).
\end{itemize}

\subsection{The minimal Galois structure that is needed}

For the dual nerve to fully activate, the KB must have relational structure
that creates $H_1$ cycles in $\Ntruth$.
This requires at minimum:
\begin{itemize}
\item A connected knowledge graph on $\geq 6$ entities.
\item At least one non-trivial relation cycle (entity $A$ relates to $B$
  relates to $C$ relates back to $A$ via a non-trivial path).
\item The generated text spanning multiple entities simultaneously.
\end{itemize}
In that setting, the comparison map $H_1(\Nmod) \to H_1(\Ntruth)$ becomes
non-trivial, and the Galois structure---governing how the KB's algebraic
relations project onto the model's mention structure---becomes operative.

% ─────────────────────────────────────────────────────────────────
\section{Empirical requirements for confirmation}
\label{sec:empirical}
% ─────────────────────────────────────────────────────────────────

\subsection{What has been confirmed}

\begin{itemize}
\item KB-similarity grounding defect: Cohen's $d = 1.82$, stable across
  12 bootstrap prompts on synthetic model.
\item $H_1$ phase-transition sign flip: factual $H_1$ growth exceeds
  hallucinated $H_1$ growth at training phase $\geq 0.3$ (synthetic phase sweep).
\item $F_\mathrm{depth} = 1$ universally: documented, root cause identified,
  metric retired from the detector.
\end{itemize}

\subsection{What requires trained checkpoints}

The following predictions are falsifiable on trained GPT-2 checkpoints
but not on synthetic models:

\begin{enumerate}
\item \textbf{SC discrimination:} $\mathrm{SC}_\mathrm{factual} \gg 1/d$ at attention sinks;
  $\mathrm{SC}_\mathrm{hallucinated} \approx 1/d$.
\item \textbf{$H_1$ direction reversal:} On trained GPT-2, factual spans should
  have higher $H_1(\alpha\text{-complex})$ than hallucinated spans,
  reversing the inversion seen on synthetic models.
\item \textbf{Coboundary gap:} $\|\rho_R - \rho_L\|_F$ should be consistently
  lower for factual windows than hallucinated windows on trained model.
\item \textbf{$H_1(\Ntruth) > 0$ enrichment:} With Wikipedia-scale KB
  ($\geq 100$ entities, relational graph with known cycles), the topological
  component of the grounding defect should activate.
\end{enumerate}

\subsection{The correct experimental setup}

Five GPT-2 checkpoints (steps 0, 1K, 10K, 100K, 500K),
FActScore labels on Wikipedia biography text,
fixed prompts and greedy decoding across all checkpoints,
and the four measurements above per atomic fact span.

% ─────────────────────────────────────────────────────────────────
\section{Conclusion}
% ─────────────────────────────────────────────────────────────────

We set out to use the algebraic topology of long generated sequences
to detect hallucination.
After four main attempts and approximately as many conceptual reversals,
the following picture has emerged.

The correct theoretical object is the comparison map
$\phi: \Nmod \to \Ntruth$ between two simplicial complexes:
one built from the model's fact-mention structure,
one built from the ground-truth knowledge graph.
Hallucination is the failure of this map to preserve the topological
features of $\Ntruth$.

This is a clean formulation that is:
\begin{itemize}
\item \textbf{Falsifiable:} it makes specific predictions about trained
  model checkpoints that can be confirmed or refuted.
\item \textbf{Grounded in mathematics:} the dual nerve is a well-defined
  object; the comparison map is a well-typed morphism of simplicial complexes.
\item \textbf{Honest about its limitations:} the topological component
  requires multi-entity KB structure to activate; current experiments
  confirm only the KB-similarity component.
\end{itemize}

The Galois thread---the three primes, the orbit taxonomy, the descent framework---
survived the successive corrections in attenuated form.
The primes remain as architectural localizations; the orbit taxonomy is retired
as a semantic classifier but retained as a context-collapse measure;
the descent framework is the correct conceptual foundation but its implementation
via supersingularity was wrong.

The failures documented here are not incidental.
They are the main content of the paper.
Each failure was preceded by a plausible theorem that cohered with the
Seidel--Henniart--Vign\'eras framework; each was detected by a specific
diagnostic experiment; each correction moved us closer to what is now
a working, if partially activated, hallucination detector.

\begin{thebibliography}{99}
\bibitem{Seidel25}
  P.~Seidel, \emph{P-adic splittings of the quantum connection},
  arXiv:2503.00500 (2025).

\bibitem{HV17}
  G.~Henniart and M.-F. Vign\'eras,
  \emph{Representations of a p-adic group over a field of characteristic $p$},
  in Representations of Reductive Groups, Proc.\ Symp.\ Pure Math., 101 (2019),
  pp.\ 171--210.

\bibitem{Borsuk48}
  K.~Borsuk, \emph{On the imbedding of systems of compacta in simplicial complexes},
  Fund.\ Math., 35 (1948), pp.\ 217--234.

\bibitem{BaudotBennequin15}
  P.~Baudot and D.~Bennequin,
  \emph{The homological nature of entropy},
  Entropy, 17(5) (2015), pp.\ 3253--3318.

\bibitem{Kontsevich14}
  M.~Kontsevich, \emph{Summation over the Integers},
  talk at SCGP, Stony Brook (2014).

\bibitem{FActScore}
  S.~Min, K.~Krishna, X.~Lyu, M.~Lewis, W.~Yih, P.~Koh, M.~Iyyer,
  L.~Zettlemoyer, and H.~Hajishirzi,
  \emph{FActScore: Fine-grained Atomic Evaluation of Factual Precision in
  Long Form Text Generation}, EMNLP 2023.
\end{thebibliography}

\end{document}
